% ID: SEL-D0080-OB-E01
% Tanda S080.
\textbf{a)} Tomamos las filas en el orden huevos, terrones y harina en gramos, y las columnas en el orden grande y pequeño:
\[
 M=\boxed{\begin{pmatrix}2&1\\5&3\\100&80\end{pmatrix}}.
\]
\textbf{b)} Con ese mismo orden de columnas, las cantidades son
\[
 A=\begin{pmatrix}20\\30\end{pmatrix},\qquad
 B=\begin{pmatrix}30\\20\end{pmatrix}.
\]
\textbf{c)} El producto suma los recursos necesarios para cada reparto:
\[
 MA=\begin{pmatrix}70\\190\\4400\end{pmatrix},
 \qquad MB=\begin{pmatrix}80\\210\\4600\end{pmatrix}.
\]
Los recursos disponibles, en las mismas unidades, son 96 huevos, 200 terrones y 5000 g de harina. Para $A$, las tres necesidades quedan por debajo de esos límites, por lo que sí se pueden preparar 20 grandes y 30 pequeños. Para $B$, se necesitan 210 terrones y solo hay 200. Por tanto, el reparto de 30 grandes y 20 pequeños no es posible.
